% MAIN SNAPSHOT — LSPV/RMV::0950::2026-10-02T09:50:00-07:00
\documentclass[11pt]{article}
\usepackage[margin=.86in]{geometry}
\usepackage{amsmath,amssymb}
\begin{document}
\section*{LSPV-5 MAIN Snapshot 0950}
\textbf{STATUS: OPEN.} P versus NP remains open. No unrestricted SAT lower bound,
arbitrary-algorithm transfer theorem, nonrelativizing hinge, or nonalgebrizing hinge is established.

\section*{CORPUS RECONCILIATION}
\textbf{Corpus surveyed.} Current repository tree/history; AGENTS/CODEX/pnp-research governance;
root \texttt{main.tex}; TASK-001--007; verified MAIN through 0858; MATRIXCERT1;
EXCESS1/EXCESSDUAL1 and concurrent TURANEX1; RAMCOMP1; predecessor TRACECLOSE1;
BITRANK1; PRIME-LQP-1; PRIMECRT1; concurrent PHCLOSE1 and PRIMEPH1; HEC/VSC-HEC/
G-PPC-HEC/CRS-PPC/CCSPC/RCCN/Search--G; active experiments/results; current GROK/CLAUDE
and independent-review artifacts.

\textbf{Duplicate theorem resolved.} PHCLOSE1 at 09:42 and PRIMEPH1 at 09:45 independently
derive the same fixed-PH prefix-search closure under \(P=NP\). Adopt PHCLOSE1's name
\[
\boxed{\text{PH-PROV-CLOSE1}}
\]
as the canonical theorem. PRIMEPH1 is retained as a direct corollary for the project's
\(\exists\sigma\forall u\exists z\forall y\) reconstruction shape and as an independent finite
implementation control, not as a second theorem promotion.

\textbf{New verified provenance result.}
If \(R(x,w)\in PH\) is polynomially balanced, then under \(P=NP\) a valid witness is selectable
in FP by bitwise prefix-extension search. Hence any polynomial-time positional-prime/CRT/source
annotation of the selected witness is also FP. Uniqueness and lexicographic canonicalization do not
escape this closure. PRIME-FNP-CLOSE1 is the first-level special case; deterministic exact histories
remain separately closed by TRACECLOSE1.

For the current reconstruction dependency
\[
\exists\sigma\forall u\exists z\forall y\,B(x,\sigma,u,z,y),\qquad B\in P,
\]
the fixed-\(\sigma\) predicate lies in \(\Pi_3^P\) and prefix extension lies in \(\Sigma_4^P\).
Thus under \(P=NP\) the outer \(\sigma\) and any polynomial-time prime annotation are
qualitatively FP-selectable. This does not supply the quantitative exponent ledger or source match.

\textbf{Alternation boundary.}
PH-PROV-CLOSE1 closes only fixed PH. Chandra--Kozen--Stockmeyer give
\[
AP=PSPACE.
\]
Polynomially many genuine alternations can therefore leave PH. This is a syntactic escape from the
closure, not a SAT lower bound: a strategy object can be exponentially large or succinctly encoded,
and no model-preserving theorem currently forces an arbitrary static SAT solver to realize such a
PSPACE-like strategy.

\textbf{Structural update.}
TURANEX1 strengthens EXCESS1/EXCESSDUAL1. For \(U\le\mathbb F_2^E\),
\[
\epsilon(I)\ge r
\iff
\exists W\le U,\ \dim W=r,\ I\text{ independent in }\Gamma(W),
\]
where \(\Gamma(W)\) is the union-support graph. If every \(s\)-set has \(\epsilon(I)<r\),
then the generalized Hamming weight satisfies
\[
\boxed{d_r(U)\ge \binom m2-t_{s-1}(m).}
\]
At \(s=\Theta(\sqrt m)\), zero excess would force \(d_{\min}(U)=\Omega(m^{3/2})\).
This is a structural code-distance theorem, not runtime hardness.

\textbf{Conflicts/stale claims resolved.}
The provenance survival condition is strengthened from not-NP-verifiable replacement to
not-polynomially-balanced-PH-verifiable replacement. Any fixed finite audit alternation is therefore
insufficient as a generic escape from the \(P=NP\) counterfactual. Prime magnitude/factorization,
deterministic trace provenance, generic canonical binding, complex Gram-to-binary-rank transfer, and
raw local-nullity abundance remain retired. TURANEX1 does not alter the arbitrary-algorithm barrier.

\section*{Provenance closure ladder}
A provenance candidate must pass:
\[
\boxed{
\begin{array}{l}
\text{semantic quotient} \Rightarrow \text{irrelevant},\\
\text{deterministic polynomial trace functional} \Rightarrow \text{TRACECLOSE1},\\
\text{polynomially balanced PH replacement relation} \Rightarrow \text{PH-PROV-CLOSE1},\\
\text{outside PH / exact-original strategy} \Rightarrow \text{still needs solver-independent necessity}.
\end{array}}
\]
Accordingly, prime provenance is instrumentation unless the fourth class is realized by an exact
source-matched theorem.

\section*{Shortest surviving proof chain}
\[
\boxed{\text{exact source-typed SU/vMS family}}
\to
\boxed{\operatorname{rank}_{\mathbb F_2}\Phi=\omega(m)\text{ or classified exceptions}}
\to
\boxed{\text{charged acquisition surviving REACQ/TRACECLOSE/PH-PROV-CLOSE}}
\]
\[
\to
\boxed{\text{model-matched reconstruction}}
\to
\boxed{\text{nonrelativizing/nonalgebrizing hinge}}
\to
\boxed{\text{unconditional contradiction}}.
\]
Every load-bearing boxed arrow remains OPEN or CONDITIONAL.

\section*{ONE SINGLE TEAM TARGET}
For one exact reconstruction-compatible SU/vMS polynomial-bit orbit, construct the source-transcribed
binary operator \(\Phi\) and legal slice matrices and prove or falsify
\[
\boxed{\operatorname{rank}_{\mathbb F_2}\Phi=\omega(m)}
\]
using intrinsically characteristic-two certificates. Run BITRANK1 odd-minor/binary-dual diagnostics,
MATRIXCERT1, EXCESS/TURANEX1 generalized-weight profiles, and COVERBREAK1 in one pipeline.
Do not promote provenance to the proof spine unless a source-specific restriction survives
PH-PROV-CLOSE1 and deterministic replay and is independently forced on arbitrary SAT deciders.

\section*{What this does not prove}
PH-PROV-CLOSE1 is a closure theorem, not a separation. TURANEX1 is a structural theorem, not a
runtime lower bound. The outside-PH alternation zone is only a candidate modeling region. P versus NP
remains OPEN.

\begin{thebibliography}{9}
\bibitem{Stockmeyer76} L. J. Stockmeyer, The Polynomial-Time Hierarchy, Theoretical Computer Science 3(1), 1976.
\bibitem{Wrathall76} C. Wrathall, Complete Sets and the Polynomial-Time Hierarchy, Theoretical Computer Science 3(1), 1976.
\bibitem{CKS81} A. K. Chandra, D. C. Kozen, L. J. Stockmeyer, Alternation, JACM 28(1), 1981.
\bibitem{MP91} N. Megiddo, C. H. Papadimitriou, On total functions, existence theorems and computational complexity, TCS 81(2), 1991.
\bibitem{KKMP21} R. Kleinberg, O. Korten, D. Mitropolsky, C. Papadimitriou, Total Functions in the Polynomial Hierarchy, ITCS 2021.
\bibitem{FIM25} N. Fleming, D. Imrek, C. Marciot, Provably Total Functions in the Polynomial Hierarchy, CCC 2025.
\bibitem{Wei91} V. K. Wei, Generalized Hamming Weights for Linear Codes, IEEE Trans. Inf. Theory 37, 1991.
\end{thebibliography}
\end{document}
% END_OF_VERIFIED_MAIN::0950
